% =========================================================
% 3.11 Flip-Flop
% =========================================================

\section{ফ্লিপ-ফ্লপ
(\textenglish{Flip-Flop})}

ডিজিটাল সার্কিটে শুধু তথ্য প্রক্রিয়াকরণ করলেই হয় না; অনেক সময় কোনো
একটি বাইনারি মান কিছু সময়ের জন্য মনে রাখতেও হয়।

যেমন, কোনো সার্কিটে বর্তমানে \(1\) সংরক্ষিত থাকলে নতুন নির্দেশ না
পাওয়া পর্যন্ত সেই \(1\) ধরে রাখতে হবে। আবার প্রয়োজন হলে সেটিকে
\(0\)-তে পরিবর্তন করতে হবে।

এই ধরনের এক-বিট তথ্য সংরক্ষণের মৌলিক সার্কিট হলো
\textbf{\textenglish{Flip-Flop}}।

সহজভাবে—

\[
\boxed{
\text{\textenglish{Input/Control Signal}}
\rightarrow
\text{\textenglish{Flip-Flop}}
\rightarrow
\text{\textenglish{Stored 0 or 1}}
}
\]


% =========================================================
% 3.11.1 What is Flip-Flop?
% =========================================================

\subsection{ফ্লিপ-ফ্লপ কী?}

যে দ্বি-স্থিতিশীল ডিজিটাল লজিক সার্কিট দুটি স্থিতিশীল অবস্থার
যেকোনো একটিতে অবস্থান করে একটি বাইনারি বিট \(0\) অথবা \(1\)
সংরক্ষণ করতে পারে, তাকে \textbf{\textenglish{Flip-Flop}} বলে।

এখানে দ্বি-স্থিতিশীল বা \textenglish{Bistable} অর্থ হলো—
সার্কিটটির দুটি স্থিতিশীল অবস্থা আছে।

\[
Q=0
\]

অথবা,

\[
Q=1
\]

অতএব,

\[
\boxed{
1\text{টি }\text{\textenglish{Flip-Flop}}
=
1\text{ }\text{\textenglish{bit}}\text{ তথ্য সংরক্ষণ}
}
\]

\textenglish{Flip-Flop}-কে তাই ডিজিটাল সার্কিটের একটি মৌলিক
\textenglish{Memory Element} বলা হয়।


% =========================================================
% 3.11.2 Output
% =========================================================

\subsection{ফ্লিপ-ফ্লপের আউটপুট
(\textenglish{Output})}

সাধারণভাবে \textenglish{Flip-Flop}-এর প্রধান আউটপুটকে \(Q\)
দ্বারা প্রকাশ করা হয়।

অনেক \textenglish{Flip-Flop}-এ এর বিপরীত আউটপুটও থাকে—

\[
\overline{Q}
\]

সাধারণ অবস্থায়,

\[
\overline{Q}=\overline{Q}
\]

অর্থাৎ—

\[
Q=1
\Rightarrow
\overline{Q}=0
\]

এবং,

\[
Q=0
\Rightarrow
\overline{Q}=1
\]

\begin{center}
\renewcommand{\arraystretch}{1.2}
\begin{tabular}{|c|c|c|}
\hline
\rowcolor{ICTLightBlue}
\(\boldsymbol{Q}\)
&
\(\boldsymbol{\overline{Q}}\)
&
\textbf{সংরক্ষিত বিট}
\\
\hline
0 & 1 & 0 \\
\hline
1 & 0 & 1 \\
\hline
\end{tabular}
\end{center}

\par\begingroup

\textbf{মনে রাখবে:}
\(Q\) এবং \(\overline{Q}\) সাধারণত পরস্পরের বিপরীত মান প্রকাশ করে।
\par\endgroup



% =========================================================
% 3.11.3 Why Flip-Flop is used
% =========================================================

\subsection{ফ্লিপ-ফ্লপ কেন ব্যবহার করা হয়?}

কম্পিউটার ও অন্যান্য ডিজিটাল ডিভাইসে অনেক কাজের জন্য পূর্ববর্তী তথ্য
বা অবস্থা মনে রাখা প্রয়োজন। \textenglish{Flip-Flop} এই তথ্য বা অবস্থা
সংরক্ষণে ব্যবহৃত হয়।

ফ্লিপ-ফ্লপ ব্যবহারের প্রধান কারণগুলো হলো—

\begin{enumerate}[leftmargin=*,itemsep=2pt]

    \item \textbf{বাইনারি তথ্য সংরক্ষণ:}
    একটি \(0\) অথবা \(1\) ধরে রাখতে।

    \item \textbf{পূর্ববর্তী অবস্থা মনে রাখা:}
    সার্কিটের বর্তমান বা পূর্বের \textenglish{state} সংরক্ষণ করতে।

    \item \textbf{\textenglish{Counter} তৈরি:}
    একাধিক \textenglish{Flip-Flop} ব্যবহার করে
    \textenglish{pulse} বা \textenglish{event} গণনা করতে।

    \item \textbf{\textenglish{Register} তৈরি:}
    একাধিক বিট একসঙ্গে সংরক্ষণ করতে।

    \item \textbf{\textenglish{Shift Register} তৈরি:}
    বিট এক অবস্থান থেকে অন্য অবস্থানে স্থানান্তর করতে।

    \item \textbf{\textenglish{Clock}-এর সঙ্গে কাজ:}
    নির্দিষ্ট সময়ে তথ্য গ্রহণ ও পরিবর্তন করতে।

    \item \textbf{\textenglish{Control Circuit}:}
    বিভিন্ন ডিজিটাল অবস্থা নিয়ন্ত্রণ করতে।

\end{enumerate}


% =========================================================
% 3.11.4 How Flip-Flop stores information
% =========================================================

\subsection{ফ্লিপ-ফ্লপ কীভাবে তথ্য সংরক্ষণ করে?}

\textenglish{Flip-Flop}-এর ভেতরে এমনভাবে
\textenglish{feedback} ব্যবস্থা তৈরি করা হয়, যাতে একবার কোনো অবস্থা
তৈরি হলে নতুন উপযুক্ত ইনপুট না আসা পর্যন্ত সেই অবস্থা ধরে রাখা যায়।

ধরা যাক,

\[
Q=1
\]

অর্থাৎ \textenglish{Flip-Flop}-এ বর্তমানে \(1\) সংরক্ষিত আছে।

নতুন কোনো পরিবর্তনকারী সংকেত না এলে—

\[
Q=1
\]

অবস্থাই বজায় থাকবে।

আবার উপযুক্ত ইনপুট দিলে অবস্থাটি পরিবর্তিত হয়ে—

\[
Q=0
\]

হতে পারে।

অর্থাৎ \textenglish{Flip-Flop} শুধু আউটপুট তৈরি করে না;
এটি আগের আউটপুট বা অবস্থা মনে রাখতে পারে।

\subsubsection*{এক বিট তথ্য সংরক্ষণের ধারণা}

\[
Q=0
\quad\Longrightarrow\quad
0\text{ সংরক্ষিত}
\]

\[
Q=1
\quad\Longrightarrow\quad
1\text{ সংরক্ষিত}
\]

\par\begingroup

\textenglish{Flip-Flop}-এর \textenglish{feedback} ব্যবস্থা
বর্তমান অবস্থা ধরে রাখতে সাহায্য করে। তবে বিদ্যুৎ সরবরাহ বন্ধ হলে
সাধারণ \textenglish{Flip-Flop}-এর সংরক্ষিত তথ্য হারিয়ে যায়।
তাই এটি একটি \textenglish{volatile storage element}।
\par\endgroup



% =========================================================
% 3.11.5 Role of Clock
% =========================================================

\subsection{ক্লকের ভূমিকা
(\textenglish{Role of Clock})}

অনেক \textenglish{Flip-Flop}-এর সঙ্গে একটি
\textenglish{Clock (CLK)} ইনপুট থাকে।

\textenglish{Clock} হলো নির্দিষ্ট সময় পরপর পরিবর্তিত হওয়া একটি
ডিজিটাল সংকেত, যা \textenglish{Flip-Flop}-কে নির্দেশ করে
কখন ইনপুট গ্রহণ বা অবস্থা পরিবর্তন করতে হবে।

\textenglish{Clock signal} সাধারণত—

\[
0\rightarrow1\rightarrow0\rightarrow1\rightarrow\cdots
\]

এভাবে পরিবর্তিত হয়।

কিছু \textenglish{Flip-Flop} ক্লকের—

\[
0\rightarrow1
\]

পরিবর্তনের সময় কাজ করে। একে
\textenglish{Positive Edge} বলা হয়।

আবার কিছু \textenglish{Flip-Flop}—

\[
1\rightarrow0
\]

পরিবর্তনের সময় কাজ করে। একে
\textenglish{Negative Edge} বলা হয়।

সহজভাবে—

\[
\boxed{
\text{\textenglish{Clock}}
=
\text{কখন }\text{\textenglish{state}}
\text{ পরিবর্তন হবে তার সময়-নিয়ন্ত্রণ সংকেত}
}
\]


% =========================================================
% 3.11.6 Characteristics
% =========================================================

\subsection{ফ্লিপ-ফ্লপের বৈশিষ্ট্য}

\begin{enumerate}[leftmargin=*,itemsep=2pt]

    \item একটি \textenglish{Flip-Flop} সাধারণত
    \(1\) \textenglish{bit} তথ্য সংরক্ষণ করে।

    \item এর দুটি স্থিতিশীল অবস্থা—\(0\) ও \(1\)।

    \item পূর্ববর্তী \textenglish{state} ধরে রাখতে পারে।

    \item সাধারণত \(Q\) এবং \(\overline{Q}\)
    আউটপুট ব্যবহৃত হয়।

    \item বিভিন্ন \textenglish{Flip-Flop} ইনপুটের নিয়ম অনুযায়ী
    অবস্থা পরিবর্তন করে।

    \item \textenglish{Clock-controlled digital circuit}-এ
    ব্যাপকভাবে ব্যবহৃত হয়।

    \item বিদ্যুৎ সরবরাহ বন্ধ হলে সাধারণ
    \textenglish{Flip-Flop}-এর সংরক্ষিত অবস্থা হারিয়ে যায়।

\end{enumerate}


% =========================================================
% 3.11.7 History
% =========================================================

\subsection{ফ্লিপ-ফ্লপের সংক্ষিপ্ত ইতিহাস}

ডিজিটাল তথ্য সংরক্ষণের ধারণা বিকাশের সঙ্গে
\textenglish{Flip-Flop}-এর ইতিহাস জড়িত।

১৯১৮ সালে \textenglish{William Eccles} ও
\textenglish{F. W. Jordan} দুটি \textenglish{vacuum tube} ব্যবহার করে
একটি দ্বি-স্থিতিশীল সার্কিট তৈরি করেন।
এটি আধুনিক \textenglish{Flip-Flop}-এর প্রাথমিক রূপগুলোর একটি।

পরবর্তী সময়ে \textenglish{vacuum tube}-এর পরিবর্তে
\textenglish{transistor} এবং পরে
\textenglish{Integrated Circuit (IC)} প্রযুক্তির সাহায্যে
\textenglish{Flip-Flop} আরও ছোট, দ্রুত ও কার্যকর হয়।

বর্তমানে \textenglish{processor}, \textenglish{register},
\textenglish{counter} এবং বিভিন্ন
\textenglish{sequential digital circuit}-এ
বিপুল সংখ্যক \textenglish{Flip-Flop} ব্যবহৃত হয়।


% =========================================================
% 3.11.8 Applications
% =========================================================

\subsection{ফ্লিপ-ফ্লপের ব্যবহার}

\begin{enumerate}[leftmargin=*,itemsep=2pt]

    \item \textbf{\textenglish{Register}:}
    একাধিক বিট সাময়িকভাবে সংরক্ষণ করতে।

    \item \textbf{\textenglish{Counter}:}
    \textenglish{clock pulse} বা
    \textenglish{event} গণনা করতে।

    \item \textbf{\textenglish{Shift Register}:}
    বাইনারি তথ্য এক অবস্থান থেকে অন্য অবস্থানে স্থানান্তর করতে।

    \item \textbf{\textenglish{Memory Circuit}:}
    ডিজিটাল তথ্য সাময়িকভাবে সংরক্ষণ করতে।

    \item \textbf{\textenglish{Frequency Divider}:}
    \textenglish{clock frequency} ভাগ বা কমাতে।

    \item \textbf{\textenglish{Control Circuit}:}
    ডিজিটাল সিস্টেমের বিভিন্ন অবস্থা নিয়ন্ত্রণ করতে।

    \item \textbf{\textenglish{Synchronizing Circuit}:}
    \textenglish{clock}-এর সঙ্গে সংকেত সমন্বয় করতে।

\end{enumerate}


% =========================================================
% 3.11.9 Types
% =========================================================

\subsection{ফ্লিপ-ফ্লপের প্রকারভেদ}

ইনপুট ও অবস্থা পরিবর্তনের পদ্ধতি অনুযায়ী
\textenglish{Flip-Flop} প্রধানত চার প্রকার—

\begin{enumerate}[leftmargin=*,itemsep=3pt]

    \item \textenglish{SR Flip-Flop (Set-Reset)}
    \item \textenglish{JK Flip-Flop}
    \item \textenglish{D Flip-Flop (Data/Delay)}
    \item \textenglish{T Flip-Flop (Toggle)}

\end{enumerate}


% =========================================================
% SR
% =========================================================

\subsubsection*{\textenglish{SR Flip-Flop (Set-Reset)}}

এর প্রধান দুটি ইনপুট—

\[
S=\text{\textenglish{Set}}
\]

এবং,

\[
R=\text{\textenglish{Reset}}
\]

\textenglish{Set} ইনপুটের মাধ্যমে আউটপুটকে \(1\)-এ
এবং \textenglish{Reset} ইনপুটের মাধ্যমে আউটপুটকে \(0\)-তে
নেওয়া যায়।

এটি \textenglish{Flip-Flop}-এর একটি মৌলিক রূপ এবং সাধারণ
অবস্থা সংরক্ষণ ও নিয়ন্ত্রণের কাজে ব্যবহৃত হয়।


% =========================================================
% JK
% =========================================================

\subsubsection*{\textenglish{JK Flip-Flop}}

\textenglish{JK Flip-Flop} হলো
\textenglish{SR Flip-Flop}-এর উন্নত রূপ।

এর বিশেষ বৈশিষ্ট্য হলো—

\[
J=K=1
\]

হলে আউটপুট \textenglish{Toggle} করে।

অর্থাৎ,

\[
0\rightarrow1
\]

অথবা,

\[
1\rightarrow0
\]

হয়।

\textenglish{Counter} এবং বিভিন্ন
\textenglish{toggle circuit}-এ এটি ব্যবহৃত হয়।


% =========================================================
% D
% =========================================================

\subsubsection*{\textenglish{D Flip-Flop (Data/Delay)}}

\textenglish{D Flip-Flop}-এর প্রধান ইনপুট হলো—

\[
D
\]

ক্লকের নির্দিষ্ট সক্রিয় মুহূর্তে \(D\)-এর মান
আউটপুট \(Q\)-তে সংরক্ষিত হয়।

সহজভাবে—

\[
\boxed{
Q_{\text{\textenglish{next}}}=D
}
\]

\textenglish{Register}, \textenglish{data storage} এবং
\textenglish{synchronization}-এর কাজে
\textenglish{D Flip-Flop} ব্যাপকভাবে ব্যবহৃত হয়।


% =========================================================
% T
% =========================================================

\subsubsection*{\textenglish{T Flip-Flop (Toggle)}}

\textenglish{T Flip-Flop}-এর প্রধান ইনপুট—

\[
T
\]

যদি,

\[
T=0
\]

হয়, তাহলে আউটপুট আগের অবস্থায় থাকে।

আর,

\[
T=1
\]

হলে আউটপুট \textenglish{Toggle} করে।

অর্থাৎ,

\[
0\leftrightarrow1
\]

এই কারণে \textenglish{T Flip-Flop}
\textenglish{Counter} ও
\textenglish{Frequency Divider}-এ বিশেষভাবে উপযোগী।


% =========================================================
% Comparison Table
% =========================================================

\subsubsection*{এক নজরে চার ধরনের ফ্লিপ-ফ্লপ}

\begin{center}

\renewcommand{\arraystretch}{1.25}

\begin{tabularx}
{\textwidth}
{|L{0.16\textwidth}|L{0.18\textwidth}|Y|Y|}

\hline

\rowcolor{ICTLightBlue}

\textbf{\textenglish{Flip-Flop}}
&
\textbf{প্রধান ইনপুট}
&
\textbf{মূল কাজ}
&
\textbf{সাধারণ ব্যবহার}
\\
\hline

\textenglish{SR}
&
\(S,R\)
&
\textenglish{Set} ও \textenglish{Reset}
&
\textenglish{State Control}
\\
\hline

\textenglish{JK}
&
\(J,K\)
&
\textenglish{Set, Reset} ও \textenglish{Toggle}
&
\textenglish{Counter}
\\
\hline

\textenglish{D}
&
\(D\)
&
\textenglish{Data} সংরক্ষণ
&
\textenglish{Register, Storage}
\\
\hline

\textenglish{T}
&
\(T\)
&
\textenglish{Toggle}
&
\textenglish{Counter, Frequency Divider}
\\
\hline

\end{tabularx}

\end{center}


% =========================================================
% Memory Trick
% =========================================================

\par\begingroup


\textbf{মনে রাখার কৌশল:}

\[
\boxed{
\text{\textenglish{SR}}
=
\text{\textenglish{Set-Reset}}
}
\]

\[
\boxed{
\text{\textenglish{JK}}
=
\text{\textenglish{Improved SR + Toggle}}
}
\]

\[
\boxed{
\text{\textenglish{D}}
=
\text{\textenglish{Data Store}}
}
\]

\[
\boxed{
\text{\textenglish{T}}
=
\text{\textenglish{Toggle}}
}
\]

\par\endgroup

\subsection{Latch, সক্রিয় ক্লক ও অবস্থা পরিবর্তন}
\textbf{ল্যাচ} নিয়ন্ত্রণের একটি স্তর সক্রিয় থাকা পর্যন্ত ইনপুটের পরিবর্তন গ্রহণ করতে পারে। \textbf{এজ-ট্রিগার্ড ফ্লিপ-ফ্লপ} ক্লকের নির্দিষ্ট ওঠা বা নামার মুহূর্তে নতুন মান গ্রহণ করে। দুইটি cross-coupled NOR গেটের সাধারণ SR latch-এ ক্লক নেই; তাকে ক্লক-নিয়ন্ত্রিত D flip-flop-এর সঙ্গে এক করে দেখানো যাবে না।

\smallhead{NOR দিয়ে active-high SR latch}
\begin{center}
\begin{minipage}{\linewidth}\centering\begin{chThreeDiagram}\begin{tikzpicture}[>=Stealth,line width=.7pt]
\node[nor gate US,chgate] (g1) at (1,1) {};
\node[nor gate US,chgate] (g2) at (1,-1) {};
\draw (g1.input 1)--++(-1,0) node[left] {$R$};
\draw (g2.input 2)--++(-1,0) node[left] {$S$};
\draw (g1.output)--(3,1) node[right] {$Q$};
\draw (g2.output)--(3,-1) node[right] {$\overline Q$};
\draw (2.3,1) to[out=-90,in=0] (2,-.15) --(-.4,-.15)|-(g2.input 1);
\draw (2.6,-1) to[out=90,in=0] (2,.15) --(-.7,.15)|-(g1.input 2);
\fill (2.3,1) circle(1.5pt);\fill (2.6,-1) circle(1.5pt);
\end{tikzpicture}\end{chThreeDiagram}
\captionof{figure}{এক NOR-এর আউটপুট অন্য NOR-এর ইনপুটে ফিরে এসে অবস্থা ধরে রাখে।}

\end{minipage}
\end{center}
$S=1,R=0$-তে Q সেট হয়; $S=0,R=1$-তে রিসেট। $S=R=0$-তে feedback আগের মান ধরে রাখে। $S=R=1$-এ দুই আউটপুটই ০ হয়; এই নিষিদ্ধ অবস্থায় পরস্পরের পূরক থাকার নিয়ম ভেঙে যায়। NAND latch-এ active-low $\overline S,\overline R$ ব্যবহৃত হয় এবং নিষিদ্ধ ইনপুট 00; দুটি সারণি একই নয়।

\smallhead{ক্লক-নিয়ন্ত্রিত ফ্লিপ-ফ্লপের কার্যসারণি}
সারণির $Q$ বর্তমান মান, $Q^+$ পরের মান। SR-তে active-high ইনপুট ধরা হয়েছে; নিচের নিয়মগুলো প্রযোজ্য সক্রিয় ক্লক মুহূর্তে এবং asynchronous preset/clear নিষ্ক্রিয় অবস্থায়।
\begin{center}\begin{tabularx}{\linewidth}{|l|c|c|Y|}\hline
ধরন & ইনপুট & $Q^+$ & কাজ\\\hline
SR & 00 & $Q$ & অপরিবর্তিত\\\hline
SR & 01 & 0 & R সক্রিয়, reset\\\hline
SR & 10 & 1 & S সক্রিয়, set\\\hline
SR & 11 & --- & নিষিদ্ধ\\\hline
JK & 00 & $Q$ & অপরিবর্তিত\\\hline
JK & 01 & 0 & reset\\\hline
JK & 10 & 1 & set\\\hline
JK & 11 & $\overline Q$ & toggle\\\hline
D & 0/1 & $D$ & ইনপুট বিট সংরক্ষণ\\\hline
T & 0 & $Q$ & অপরিবর্তিত\\\hline
T & 1 & $\overline Q$ & toggle\\\hline
\end{tabularx}\end{center}

\chThreeExample{১}{D flip-flop-এ তিন ক্লক}
শুরুতে $Q=0$। সক্রিয় ওঠা প্রান্তে D-এর মান যথাক্রমে ১, ০, ১।
\chThreeAnswer ক্লকের পরে Q যথাক্রমে ১, ০, ১। দুই সক্রিয় প্রান্তের মাঝখানে D বদলালেও Q নিজের থেকে বদলাবে না।

\chThreeExample{২}{JK ও T-তে toggle}
JK-তে $J=K=1$ এবং শুরুতে $Q=0$। চার সক্রিয় ক্লকের পর Q কত?
\chThreeAnswer $0\to1\to0\to1\to0$; Q=০। T=১ করে একই ক্রম পাওয়া যায়। প্রতি দুই ইনপুট ক্লকে একটি পূর্ণ output cycle; আদর্শভাবে $f_Q=f_{CLK}/2$।

\chThreeExample{৩}{এক বিট কেন, দুই বিট নয়?}
Q ও $\overline Q$ দুটি output থাকলেও একটি flip-flop এক বিট রাখে।
\chThreeAnswer দুটি output স্বাধীন নয়; Q-এর মান নির্ধারিত হলেই অন্যটি নির্ধারিত। বৈধ স্থিতিশীল অবস্থা মাত্র দুটি, তাই ধারণক্ষমতা এক বিট।

